% !TeX program = xelatex
% 从工程根目录运行：latexmk examples/code-styles.tex
\documentclass[10pt,letterpaper]{book}
% 原生 XeLaTeX 排版配置。字体不依赖 Windows，也不需要 shell escape。
\usepackage{iftex}
\ifXeTeX\else
  \PackageError{designpractice}{Please compile with XeLaTeX}{Use latexmk main.tex.}
\fi
\usepackage{amsmath,amssymb}
\usepackage{unicode-math}
\defaultfontfeatures{}
\usepackage[fontset=fandol]{ctex}
\setmainfont{lmroman10-regular.otf}[
  BoldFont=lmroman10-bold.otf,
  ItalicFont=lmroman10-italic.otf,
  BoldItalicFont=lmroman10-bolditalic.otf,
  Ligatures=TeX
]
\setsansfont{lmsans10-regular.otf}[
  BoldFont=lmsans10-bold.otf,
  ItalicFont=lmsans10-oblique.otf,
  BoldItalicFont=lmsans10-boldoblique.otf,
  Ligatures=TeX
]
% 接近原稿的等宽字体；特殊符号由 code-style.tex 中的 DejaVu 回退补齐。
\setmonofont{lmmono10-regular.otf}[
  BoldFont=lmmonolt10-bold.otf,
  ItalicFont=lmmono10-italic.otf,
  BoldItalicFont=lmmonolt10-boldoblique.otf,
  Ligatures=NoCommon
]
\setCJKfamilyfont{bookmono}{FandolFang-Regular.otf}[
  BoldFont=FandolHei-Bold.otf,
  ItalicFont=FandolKai-Regular.otf,
  BoldItalicFont=FandolHei-Bold.otf
]
\renewcommand{\CJKttdefault}{bookmono}

% 沿用原书开本和页边距；正文允许自然分页。
\usepackage[left=2cm,right=2cm,top=3.5cm,bottom=2.5cm]{geometry}
\usepackage{parskip}
\usepackage{microtype}
\UseMicrotypeSet[protrusion]{basicmath}
\raggedbottom
\setlength{\emergencystretch}{3em}
\setcounter{secnumdepth}{5}
\setcounter{tocdepth}{1}

\usepackage{xcolor}
\usepackage{upquote}
% 原生 LaTeX 语法高亮；修改代码后随 XeLaTeX 编译自动更新，无预生成缓存。
\usepackage{listings}
\usepackage{xeCJK-listings}

\definecolor{CodeKeyword}{HTML}{1F4E79}
\definecolor{CodeType}{HTML}{68428A}
\definecolor{CodeString}{HTML}{9A4B00}
\definecolor{CodeComment}{HTML}{446845}
\definecolor{CodeAdded}{HTML}{24683D}
\definecolor{CodeRemoved}{HTML}{A02C35}
\definecolor{CodeBackground}{HTML}{FAFAF8}
\definecolor{CodeRule}{HTML}{D7DDE3}

% 程序使用 Latin Modern Mono；纯文本/目录树使用覆盖特殊符号的 DejaVu。
\newfontfamily\CodeSymbolFont{DejaVuSansMono.ttf}[
  Path=fonts/,Scale=0.872,Ligatures=NoCommon,
  BoldFont=DejaVuSansMono-Bold.ttf,
  ItalicFont=DejaVuSansMono-Oblique.ttf,
  BoldItalicFont=DejaVuSansMono-BoldOblique.ttf
]
\lstdefinestyle{bookcode}{
  basicstyle=\small\ttfamily,
  keywordstyle=\color{CodeKeyword}\bfseries,
  keywordstyle=[2]\color{CodeType},
  keywordstyle=[3]\color{CodeKeyword},
  commentstyle=\color{CodeComment},
  stringstyle=\color{CodeString},
  identifierstyle=\color{black},
  backgroundcolor=\color{CodeBackground},
  rulecolor=\color{CodeRule},
  frame=tb,
  framerule=0.4pt,
  framesep=2mm,
  aboveskip=0.8\baselineskip,
  belowskip=0.8\baselineskip,
  columns=fullflexible,
  keepspaces=true,
  showstringspaces=false,
  showspaces=false,
  showtabs=false,
  tabsize=4,
  upquote=true,
  breaklines=true,
  breakatwhitespace=false,
  breakautoindent=true,
  breakindent=1em,
  numbers=none,
  mathescape=false,
  texcl=false,
  escapeinside={}
}

\lstdefinelanguage{BookText}{basicstyle=\small\ttfamily\CodeSymbolFont}
\lstdefinelanguage{BookVerilog}[]{Verilog}{
  morekeywords=[2]{signed,unsigned,integer,reg,wire},
  morekeywords=[3]{`define,`include,`ifdef,`ifndef,`endif,`timescale}
}
\lstdefinelanguage{BookSystemVerilog}[]{BookVerilog}{
  morekeywords={always_comb,always_ff,always_latch,logic,bit,byte,int,longint,
    shortint,struct,typedef,enum,package,endpackage,import,interface,endinterface,
    modport,unique,priority,assert,property,endproperty,sequence,endsequence}
}
\lstdefinelanguage{BookScala}[]{Scala}{
  morekeywords=[2]{UInt,Bool,Bits,SInt,Vec,Bundle,Module,Component,Input,Output,
    Reg,RegInit,RegNext,Wire,Mux,Cat,Fill,U,S,B,True,False},
  morekeywords={when,elsewhen,otherwise,switch,is}
}
\lstdefinelanguage{BookC}[]{C}{
  morekeywords=[2]{uint8_t,uint16_t,uint32_t,uint64_t,int8_t,int16_t,int32_t,
    int64_t,uintptr_t,size_t,bool,u8,u16,u32,u64,NULL},
  morekeywords=[3]{__init,__initdata,__iomem,__asm__,__volatile__}
}
\lstdefinelanguage{BookShell}[]{bash}{
  morekeywords=[2]{make,git,cp,mkdir,rm,ls,sudo,chmod,tar,grep,find,ifconfig,ping,
    setenv,load,tftpboot,bootelf,nproc,dc_shell,verilator,strip,gcc,export,
    vcs,simv,verdi,dve,java,pwd,sbt,gtkwave,g,clear,
    loongarch32r-linux-gnusf-strip,loongarch32r-linux-gnusf-gcc,
    loongarch32r-linux-gnusf-objdump,loongarch32r-linux-gnusf-objcopy}
}
\lstdefinelanguage{BookMake}[gnu]{make}{
  morekeywords=[2]{CONFIG_LS_SOC,CONFIG_BX_SOC,CONFIG_MEGA_SOC_FPGA}
}
\lstdefinelanguage{BookTcl}[]{tcl}{
  morekeywords=[2]{set_app_var,read_verilog,read_file,analyze,elaborate,link,
    uniquify,check_design,compile,compile_ultra,check_timing,report_timing,
    set_synthesis_strategy,create_clock,get_ports,get_nets,get_cells,get_clocks,
    all_inputs,all_outputs,set_input_delay,set_output_delay,set_clock_latency,
    set_clock_uncertainty,set_clock_transition,set_propagated_clock,
    set_max_transition,set_max_fanout,set_max_capacitance,set_max_area,
    set_false_path,set_multicycle_path,write,write_file,write_sdf}
}
\lstdefinelanguage{BookDTS}{
  sensitive=true,
  alsoletter={-\#},
  morekeywords={compatible,reg,interrupts,interrupt-parent,status,ranges,
    device_type,clock-frequency,clocks,clock-names,phandle,model,bootargs,
    gpio-controller,interrupt-controller,\#address-cells,\#size-cells,
    \#interrupt-cells,\#gpio-cells},
  morekeywords=[2]{GIC_SPI,IRQ_TYPE_LEVEL_HIGH,IRQ_TYPE_EDGE_RISING},
  morecomment=[l]{//},morecomment=[s]{/*}{*/},morestring=[b]"
}
\lstdefinelanguage{BookLoongArch}{
  sensitive=true,alsoletter={.},alsoother={$},
  morekeywords={add.w,addi.w,sub.w,lu12i.w,lu32i.d,lu52i.d,li.w,li.d,la.local,
    la.global,la.abs,la.pcrel,or,ori,and,andi,xor,xori,nor,slt,sltu,slti,sltui,
    sll.w,slli.w,srl.w,srli.w,sra.w,srai.w,mul.w,mulh.w,mulh.wu,div.w,div.wu,
    mod.w,mod.wu,ld.w,ld.b,ld.bu,ld.h,ld.hu,st.w,st.b,st.h,ll.w,sc.w,
    jirl,bl,b,beq,bne,beqz,bnez,bge,bgeu,blt,bltu,blez,bgtz,bltz,bgez,
    csrrd,csrwr,csrxchg,iocsrrd.w,iocsrwr.w,ertn,syscall,break,tlbsrch,tlbrd,
    tlbwr,tlbfill,invtlb,cacop,dbar,ibar,nop,move,ret},
  morekeywords=[3]{.text,.data,.bss,.section,.align,.globl,.global,.type,.size,
    .macro,.endm,.word,.long,.byte,.space,.equ,.set,.rept,.endr},
  morekeywords=[2]{zero,ra,tp,sp,a0,a1,a2,a3,a4,a5,a6,a7,t0,t1,t2,t3,t4,t5,
    t6,t7,t8,r0,r1,r2,r3,r4,r5,r6,r7,r8,r9,r10,fp,s0,s1,s2,s3,s4,s5,s6,s7,s8},
  morecomment=[l]{\#},morecomment=[l]{//},morecomment=[s]{/*}{*/},
  morestring=[b]"
}
\lstdefinelanguage{BookPseudo}[]{Python}{
  morekeywords={true,false,elseif,then,end,foreach,repeat,until}
}
\lstdefinelanguage{BookConfig}{
  sensitive=true,identifierstyle=\color{CodeType},
  morekeywords={config,menuconfig,menu,endmenu,choice,endchoice,bool,tristate,
    string,int,hex,prompt,default,depends,on,select,imply,if,endif,source,help},
  morekeywords=[2]{y,n,m},morecomment=[l]{\#},morestring=[b]"
}
\lstdefinelanguage{BookLinker}{
  sensitive=true,
  morekeywords={OUTPUT_FORMAT,OUTPUT_ARCH,ENTRY,MEMORY,SECTIONS,INCLUDE,
    ORIGIN,LENGTH,ALIGN,KEEP,PROVIDE,AT,LOADADDR,ADDR,SIZEOF},
  morecomment=[s]{/*}{*/},morecomment=[l]{\#},morestring=[b]"
}
\lstdefinelanguage{BookDiff}{
  sensitive=true,
  morecomment=[l][\color{CodeAdded}]{+},
  morecomment=[l][\color{CodeRemoved}]{-},
  morecomment=[l][\color{CodeType}]{@@},
  morecomment=[l][\color{CodeKeyword}\bfseries]{diff},
  morecomment=[l][\color{CodeKeyword}]{index}
}

% 例如：\begin{CodeBlock}[language=BookVerilog] ... \end{CodeBlock}
\lstnewenvironment{CodeBlock}[1][language=BookText]
  {\lstset{style=bookcode,#1}}
  {}


\usepackage{longtable,booktabs,array,calc}
\usepackage{etoolbox}
\AtBeginEnvironment{longtable}{\small\setlength{\tabcolsep}{4pt}}
\makeatletter
\patchcmd\longtable{\par}{\if@noskipsec\mbox{}\fi\par}{}{}
\makeatother
\usepackage{footnotehyper}
\makesavenoteenv{longtable}
\usepackage{graphicx}
\makeatletter
\def\maxwidth{\ifdim\Gin@nat@width>\linewidth\linewidth\else\Gin@nat@width\fi}
\def\maxheight{\ifdim\Gin@nat@height>\textheight\textheight\else\Gin@nat@height\fi}
\setkeys{Gin}{width=\maxwidth,height=\maxheight,keepaspectratio}
\def\fps@figure{htbp}
\makeatother
\providecommand{\tightlist}{\setlength{\itemsep}{0pt}\setlength{\parskip}{0pt}}
\usepackage{amsthm}
\usepackage[unicode]{hyperref}
\usepackage{bookmark}
\usepackage{xurl}
\urlstyle{same}

\newfontfamily\PreviousCodeFont{DejaVuSansMono.ttf}[
  Path=fonts/,Scale=0.9,Ligatures=NoCommon,
  BoldFont=DejaVuSansMono-Bold.ttf,
  ItalicFont=DejaVuSansMono-Oblique.ttf,
  BoldItalicFont=DejaVuSansMono-BoldOblique.ttf
]
\begin{document}
\pagestyle{plain}
\section*{代码字体与语法高亮示例}

下面两段使用相同代码和配色，可以直接比较字体。新版采用接近原稿的
Latin Modern Mono；中文注释使用 Fandol 仿宋风格。纯文本和目录树继续使用
DejaVu，以完整显示线条和特殊空格。

\subsection*{新版程序字体：Latin Modern Mono}
\begin{CodeBlock}[language=BookVerilog]
// 时序逻辑：低电平复位，保留 0/O、1/l、下划线和引号的区别
module counter(input wire clk, reset_n, output reg [7:0] count);
  always @(posedge clk) begin
    if (!reset_n) count <= 8'h00;
    else count <= count + 1'b1;
  end
  initial $display("counter is ready");
endmodule
\end{CodeBlock}

\subsection*{迁移初版字体：DejaVu Sans Mono（同配色对照）}
\begin{CodeBlock}[language=BookVerilog,basicstyle=\small\ttfamily\PreviousCodeFont]
// 时序逻辑：低电平复位，保留 0/O、1/l、下划线和引号的区别
module counter(input wire clk, reset_n, output reg [7:0] count);
  always @(posedge clk) begin
    if (!reset_n) count <= 8'h00;
    else count <= count + 1'b1;
  end
  initial $display("counter is ready");
endmodule
\end{CodeBlock}

\subsection*{目录树与特殊字符}
\begin{CodeBlock}[language=BookText]
chip
├── rtl  处理器设计
│   ├── core.sv
│   └── soc.sv
└── sim  仿真环境
ASCII: 0 O 1 l I  _ - --  ' "  $ # % & < >
\end{CodeBlock}

\clearpage
\section*{通用语言与硬件设计}
\subsection*{C：关键字、注释、字符串}
\begin{CodeBlock}[language=BookC]
#include <stdint.h>
// 中文注释应完整显示，并与英文注释保持同色。
uint32_t read_value(volatile uint32_t *address) {
    if (address == NULL) return 0;
    printf("value = %u\n", *address);
    return *address;
}
\end{CodeBlock}

\subsection*{Scala / Chisel / SpinalHDL}
\begin{CodeBlock}[language=BookScala]
val count = RegInit(0.U(8.W))
when (io.enable) {
  count := count + 1.U
}.otherwise {
  count := 0.U
}
\end{CodeBlock}

\subsection*{SystemVerilog}
\begin{CodeBlock}[language=BookSystemVerilog]
logic [31:0] data;
always_ff @(posedge clk) begin
  if (!reset_n) data <= '0;
  else data <= next_data;
end
\end{CodeBlock}

\subsection*{Shell}
\begin{CodeBlock}[language=BookShell]
# 编译命令保持原始内容，不由 LaTeX 执行。
export ARCH=loongarch
make -j$(nproc)
echo "Build completed"
\end{CodeBlock}

\clearpage
\section*{本书专用语言规则}
\subsection*{LoongArch 汇编}
\begin{CodeBlock}[language=BookLoongArch]
.text
.globl reset_entry
reset_entry:
    li.w    $t0, 0x1c000000  # 初始化地址
    ld.w    $a0, $t0, 0
    addi.w  $a0, $a0, 1
    st.w    $a0, $t0, 0
    ertn
\end{CodeBlock}

\subsection*{设备树 DTS}
\begin{CodeBlock}[language=BookDTS]
uart0: serial@1fe001e0 {
    compatible = "ns16550a";
    reg = <0x1fe001e0 0x8>;
    interrupts = <GIC_SPI 2 IRQ_TYPE_LEVEL_HIGH>;
    status = "okay";
};
\end{CodeBlock}

\subsection*{EDA Tcl}
\begin{CodeBlock}[language=BookTcl]
# 创建时钟并设置输入延迟。
set target_library "typical.db"
read_verilog design.v
create_clock -name clk -period 10 [get_ports clk]
set_input_delay -clock clk 3 [all_inputs]
compile -map_effort high
\end{CodeBlock}

\clearpage
\section*{配置、构建与补丁}
\subsection*{Kconfig / defconfig}
\begin{CodeBlock}[language=BookConfig]
config MEGA_SOC_FPGA
    bool "MEGA SoC for FPGA"
    default y
# CONFIG_DEBUG_INFO is not set
CONFIG_LOG_BUF_SHIFT=18
\end{CodeBlock}

\subsection*{Makefile}
\begin{CodeBlock}[language=BookMake]
CC = loongarch32r-linux-gnusf-gcc
all: kernel.o
kernel.o: kernel.c
	$(CC) -c kernel.c -o kernel.o
\end{CodeBlock}

\subsection*{链接脚本}
\begin{CodeBlock}[language=BookLinker]
OUTPUT_FORMAT(elf32-loongarch)
MEMORY {
  ram (rwxa) : ORIGIN = 0x80000000, LENGTH = 128M
}
INCLUDE "section.ld"
\end{CodeBlock}

\clearpage
\section*{伪代码与补丁}
\subsection*{伪代码}
\begin{CodeBlock}[language=BookPseudo]
# 查找命中的 TLB 条目
for entry in tlb:
    if entry.valid and entry.vpn == vpn:
        return entry.ppn
return None
\end{CodeBlock}

\subsection*{Diff：增加为绿色，删除为红色}
\begin{CodeBlock}[language=BookDiff]
diff --git a/main.c b/main.c
@@ -1,2 +1,2 @@
-return old_value;
+return new_value;
\end{CodeBlock}
\end{document}
